\section{Cancellation by one additional central field}
\label{sec:central-extension}

The hypothesis that the gauge algebra is semisimple is essential to the
no-go statement.  A central field supplies a symmetric invariant cubic
even when the simple algebra has none.  The resulting cancellation must
also satisfy the equations with central external fields.  We compute all
of these equations below.

Extend the simple algebra $\g$ by a central generator $\mathbf 1$ and
extend its invariant form by
\begin{equation}
 \kappa(X,\mathbf1)=0,\qquad
 \kappa(\mathbf1,\mathbf1)=\ell\ne0,\qquad X\in\g.
 \label{eq:central-pairing}
\end{equation}
The symmetric invariant cubic $T$ is specified by
\begin{equation}
\begin{split}
 T(X,Y,\mathbf1)&=\chi\kappa(X,Y),\qquad
 T(\mathbf1,\mathbf1,\mathbf1)=\gamma,\\
 T(X,Y,Z)&=0,\qquad T(X,\mathbf1,\mathbf1)=0,
 \qquad X,Y,Z\in\g.
\end{split}
\label{eq:central-cubic}
\end{equation}
All evaluations not displayed are fixed by symmetry and multilinearity.
Its Taylor cochain is
\begin{equation}
 j(X_0e^{p_0z},X_1e^{p_1z},X_2e^{p_2z})
   =[p_0p_1]T(X_0,X_1,X_2),\qquad p_0+p_1+p_2=0.
 \label{eq:central-j}
\end{equation}

\begin{proposition}[The complete central-field square]
\label{prop:central-square}
The cochain \eqref{eq:central-j} is classically closed.
For four external momenta with sum zero, define the already occurring
metric four-cochain by
\begin{equation}
 G_{kk}=\alt\bigl([p_0p_1][p_2p_3]
                 \kappa(X_0,X_1)\kappa(X_2,X_3)\bigr),
 \qquad X_i\in\g.
\end{equation}
The values of $\tfrac12\{j,j\}$ on every central-input sector are
\begin{equation}
\begin{array}{c|l}
 \text{number of central inputs}&\tfrac12\{j,j\}\\ \hline
 0&\displaystyle \frac{\chi^2}{8\ell}\,G_{kk}\\[2mm]
 1&0\\[1mm]
 2&\displaystyle
 \chi\left(\frac\gamma\ell-\chi\right)
 [p_0p_1][p_2p_3]\kappa(X_0,X_1),
 \quad (X_2,X_3)=(\mathbf1,\mathbf1)\\[2mm]
 3&0\\[1mm]
 4&0.
\end{array}
\label{eq:central-square-sectors}
\end{equation}
The two-central value in any other order follows by alternation.
For $\chi\ne0$, eliminating the two-central obstruction requires
\begin{equation}
 \gamma=\ell\chi.
 \label{eq:central-necessary}
\end{equation}
\end{proposition}

\begin{proof}
The only potentially nonzero classical differential of the mixed cubic
has three simple inputs and one central input.  Its value is
\begin{equation}
 -\chi\kappa([X_0,X_1],X_2)
 \bigl([p_0p_3]+[p_1p_3]+[p_2p_3]\bigr)=0.
\end{equation}
The pure central cubic has zero classical differential.  Thus $dj=0$.

The six $(2,2)$ shuffles in the bracket occur in equal pairs.  After the
factor $1/2$, their sum is the three pair partitions, with signs
$+,-,+$:
\begin{align}
 \tfrac12\{j,j\}
 &= [p_0p_1][p_2p_3]
       \kappa^{-1}\bigl(T(X_0,X_1,\cdot),T(\cdot,X_2,X_3)\bigr)\notag\\
 &\quad-[p_0p_2][p_1p_3]
       \kappa^{-1}\bigl(T(X_0,X_2,\cdot),T(\cdot,X_1,X_3)\bigr)\notag\\
 &\quad+[p_0p_3][p_1p_2]
       \kappa^{-1}\bigl(T(X_0,X_3,\cdot),T(\cdot,X_1,X_2)\bigr).
 \label{eq:central-three-pairings}
\end{align}
For four simple inputs the contracted field is central, so its inverse
pairing is $\ell^{-1}$.  The full alternation defining $G_{kk}$ is eight
times the three pairings in \eqref{eq:central-three-pairings}.  This gives
the first row of \eqref{eq:central-square-sectors}.

For $(X_0,X_1,\mathbf1,\mathbf1)$, the first pairing contributes
$\chi\gamma\kappa(X_0,X_1)/\ell$.  The other pairings each contribute
$\chi^2\kappa(X_0,X_1)$.  The exact determinant identity
\begin{equation}
 [p_0p_1][p_2p_3]-[p_0p_2][p_1p_3]
       +[p_0p_3][p_1p_2]=0
 \label{eq:central-plucker}
\end{equation}
then gives the two-central row.  The same identity kills the
four-central row, whose common color coefficient is $\gamma^2/\ell$.
The one- and three-central rows vanish from
\eqref{eq:central-cubic} before the momentum sum.

To prove necessity of \eqref{eq:central-necessary}, take
$X_0=X_1=X$ with $\kappa(X,X)\ne0$.  These two inputs and the two
central inputs commute.  Every classical boundary vanishes on them,
whereas $[p_0p_1][p_2p_3]$ is not the zero polynomial on conserved
momenta.  Bare diagrams also vanish when any external field is central.
Consequently a nonzero two-central row cannot be removed by an
order-$\hbar^2$ counterterm.
\end{proof}

\begin{theorem}[A full two-loop central cancellation]
\label{thm:central-cancellation}
Use the Killing form on $\g$.  Suppose the original graph tensors obey
the cochain identities
\begin{equation}
 F_{133}+F_{223}=\tfrac12\beta,\qquad
 F_{223}=x\beta+aG_{kk},\qquad
 \beta=d\bigl([p_0p_1]^2\kappa(X_0,[X_1,X_2])\bigr).
 \label{eq:central-graph-assumption}
\end{equation}
Then the theory for $\g\oplus\mathbb C\mathbf1$ has a solution of the
complete local quantum master equation through order $\hbar^2$ with
\begin{equation}
 \boxed{\quad
 \gamma=\ell\chi,\qquad
 \chi^2=\frac{\epsilon U\ell a}{384\pi^4},\qquad
 i_2=\frac{\epsilon U}{\pi^4}
       \left(\frac1{10240}-\frac{x}{3072}\right)
       [p_0p_1]^2\kappa(X_0,[X_1,X_2]).\quad}
 \label{eq:central-solution}
\end{equation}
Here $i_2$ vanishes if one of its arguments is central, and the
coefficients are taken over $\mathbb C$.  The signs $\epsilon$ and $U$
are the tree and loop signs of the main normalization.
\end{theorem}

\begin{proof}
A central external field kills every bare cubic graph because its
adjoint action is zero.  The bare two-loop cochain is therefore the
simple-algebra cochain, extended by zero, and
\begin{equation}
 o_2=-\frac{U}{64\pi^4}
       \left[\left(\frac{x}{48}-\frac1{160}\right)\beta
                  +\frac{a}{48}G_{kk}\right].
 \label{eq:central-bare-anomaly}
\end{equation}
The chosen central cubic cancels the two-central row in
\eqref{eq:central-square-sectors}.  The coefficient of $G_{kk}$ in
$\tfrac12\{j,j\}+\epsilon o_2$ is exactly
\begin{equation}
 \frac{\chi^2}{8\ell}-\frac{\epsilon U a}{3072\pi^4}=0.
\end{equation}
The remaining coefficient is the classical differential of the
displayed $i_2$.  Thus
$-di_2+\tfrac12\{j,j\}+\epsilon o_2=0$ on every input, which is the
full order-$\hbar^2$ equation in the stated descent convention.
The order-$\hbar$ equation was proved above.  The local one-loop
contacts containing $j$ vanish by the odd-loop theorem.
\end{proof}

The hypothesis \eqref{eq:central-graph-assumption} is an identity of
entire cochains.  A numerical contraction detecting one nonzero class
does not by itself establish this hypothesis.  The previously computed
identities in \cite{RepoOriginal} supply examples; the following one is
also recomputed directly in the present paper's certificate.

\begin{corollary}[An explicit $\mathfrak{so}(5)\oplus\mathbb C$ solution]
\label{cor:so5-central}
For $\g=\mathfrak{so}(5)$ with its Killing form,
\begin{equation}
 F_{223}=-\beta-\frac1{18}G_{kk}.
\end{equation}
Consequently
\begin{equation}
 \chi^2=-\frac{\epsilon U\ell}{6912\pi^4},\qquad
 \gamma=\ell\chi,\qquad
 i_2=\frac{13\epsilon U}{30720\pi^4}
          [p_0p_1]^2\kappa(X_0,[X_1,X_2])
 \label{eq:so5-central-solution}
\end{equation}
solves the complete local equation through order $\hbar^2$.
At $\ell=1$ and $\epsilon U=1$, the two first-order choices are
\begin{equation}
 \chi=\gamma=\frac{i}{48\sqrt3\,\pi^2}
 \qquad\text{or}\qquad
 \chi=\gamma=-\frac{i}{48\sqrt3\,\pi^2}.
\end{equation}
\end{corollary}

For completeness, the full graph identities already certified
in \cite{RepoOriginal} for the five exceptional simple algebras
give the following additional instances of
Theorem~\ref{thm:central-cancellation}:
\begin{center}
\renewcommand{\arraystretch}{1.6}
\begin{tabular}{@{}lcc@{}}
\toprule
$\g$ & $\chi^2\pi^4/(\epsilon U\ell)$
 & Coefficient of $[p_0p_1]^2f$ in $\pi^4 i_2/(\epsilon U)$\\
\midrule
$G_2$ & $-\dfrac5{36864}$ & $\dfrac{71}{368640}$\\
\addlinespace[3pt]
$F_4$ & $-\dfrac5{124416}$ & $\dfrac{41}{829440}$\\
\addlinespace[3pt]
$E_6$ & $-\dfrac1{36864}$ & $\dfrac{11}{368640}$\\
\addlinespace[3pt]
$E_7$ & $-\dfrac1{62208}$ & $\dfrac{11}{829440}$\\
\addlinespace[3pt]
$E_8$ & $-\dfrac1{115200}$ & $\dfrac1{460800}$\\
\bottomrule
\end{tabular}
\end{center}
In each row $\gamma=\ell\chi$. These are algebraic consequences
of the cited complete cochain identities. The new independent
original-graph replay described next is for $\mathfrak{so}(5)$.

\paragraph{Exact certificate.}
The script \texttt{checks/central\_extension\_exact.py} constructs the
ten matrices $E_{ij}-E_{ji}$ in $\mathfrak{so}(5)$ and checks their
complete commutator table, Jacobi identity, and Killing form
$\kappa=-6\delta$.  It contracts the six original structure tensors of
each theta graph, with no graph-engine imports, then performs all
$24$ alternations and checks all six coefficients of
$[p_0p_1]^2,[p_0p_2]^2,[p_1p_2]^2,
[p_0p_1][p_0p_2],[p_0p_1][p_1p_2],[p_0p_2][p_1p_2]$.
Each graph identity is checked on $60000$ color--momentum entries.
The central extension has respectively $10000,4000,600,40,1$
ordered color assignments with $0,1,2,3,4$ central inputs.
The full bracket certificate checks $351384$ coefficient entries;
the assembled quantum master equation checks $87846$ entries.
The arithmetic is integral after explicit denominator clearing.

This establishes a two-loop cancellation for a reductive algebra
whose simple factor is obstructed on its own.  It does not establish
an all-order quantization for that reductive algebra: higher
obstructions require the corresponding local cohomology and graph
consistency analysis.
